\documentclass[11pt,twoside]{amsart}
\usepackage{amssymb, amsmath, enumerate, palatino, hyperref}
\usepackage[normalem]{ulem}
\usepackage{fullpage}
\usepackage[T1]{fontenc}
\renewcommand{\labelitemi}{\guillemotright}
\usepackage{mathrsfs}

\usepackage{tikz}
\usepackage{tikz-cd}
\usepackage{pgfplots}
\pgfplotsset{compat=1.13}
\usetikzlibrary{positioning}
\usetikzlibrary{trees}
\usetikzlibrary{decorations.pathmorphing}    
\usetikzlibrary{decorations.pathreplacing}
\usetikzlibrary{arrows.meta,calc}
\usetikzlibrary{bending}
\usetikzlibrary{decorations.markings,shapes.geometric}
\tikzset{->-/.style={decoration={markings, mark=at position #1 with
  {\arrow{>}}},postaction={decorate}}}
\tikzset{-|-/.style={decoration={markings, mark=at position #1 with
  {\arrow{stealth}}},postaction={decorate}}}
\tikzset{movearrow/.style 2 args ={
        decoration={markings,
	  mark= at position {#1} with {\arrow{#2}} ,
        },
        postaction={decorate}
    }
}
\tikzset{<--/.style={decoration={markings, mark=at position #1 with
  {\arrow{<}}},postaction={decorate}}}
\tikzstyle{ball} = [circle,shading=ball, ball color=black,
    minimum size=1mm,inner sep=1.3pt]
\tikzstyle{bball} = [circle,shading=ball, ball color=blue,
    minimum size=1mm,inner sep=1.3pt]
\tikzstyle{miniball} = [circle,shading=ball, ball color=black,
    minimum size=1mm,inner sep=0.5pt]
\tikzstyle{redminiball} = [circle,shading=ball, ball color=red,
    minimum size=1mm,inner sep=0.5pt]
\tikzstyle{mminiball} = [circle,shading=ball, ball color=black,
    minimum size=0.6mm,inner sep=0.1pt]
\tikzstyle{mycircle} = [draw,circle, minimum size=1mm,inner sep=1.3pt]


\theoremstyle{plain}
\newtheorem{prop}{Proposition}%[section]
\newtheorem{lemma}[prop]{Lemma}
\newtheorem{thm}[prop]{Theorem}
\newtheorem{obs}[prop]{Observation}
\newtheorem{app}[prop]{Application}
\newtheorem*{MainThm}{Main Theorem}
\newtheorem{cor}[prop]{Corollary}
\newtheorem{conj}[prop]{Conjecture}
\theoremstyle{remark}
\newtheorem{rmk}[prop]{Remark}
\newtheorem{prob}{Problem}
\newtheorem{bonus}[prop]{Bonus Problem}
\newtheorem{exc}{Exercise}
\theoremstyle{definition}
\newtheorem{ex}[prop]{Example}
\theoremstyle{definition}
\newtheorem{defn}[prop]{Definition}

\newcommand{\RR}{\mathbb{R}}
\newcommand{\ZZ}{\mathbb{Z}}
\newcommand{\CC}{\mathbb{C}}
\newcommand{\NN}{\mathbb{N}}
\newcommand{\QQ}{\mathbb{Q}}
\newcommand{\PP}{\mathbb{P}}

\newcommand{\cC}{\mathscr{C}}
\newcommand{\Ob}{\operatorname{Ob}}
\newcommand{\Mor}{\operatorname{Mor}}

\newcommand{\Set}{\operatorname{Set}}
\newcommand{\FinSet}{\operatorname{FinSet}}
\newcommand{\Vect}{\operatorname{Vect}}
\newcommand{\FinVect}{\operatorname{FinVect}}
\newcommand{\Mat}{\operatorname{Mat}}
\newcommand{\Gp}{\operatorname{Gp}}
\newcommand{\FinGp}{\operatorname{FinGp}}
\newcommand{\AbGp}{\operatorname{AbGp}}
\newcommand{\Ring}{\operatorname{Ring}}
\newcommand{\CommRing}{\operatorname{CommRing}}
\newcommand{\Field}{\operatorname{Field}}
\newcommand{\Top}{\operatorname{Top}}
\newcommand{\Aut}{\operatorname{Aut}}
\newcommand{\id}{\operatorname{id}}


\newcommand{\defeq}{:=}

\title{Math 113: Discrete Structures\\ Homework 16}
%\author{} %Remove the leading % and put your name here if using this .tex file as a template for your homework.

\begin{document}
\maketitle

\noindent
\textbf{Due:} Friday, March 6 at 10pm.
\bigskip

\begin{prob}
Consider the following floor plan for a building:

\begin{center}
\begin{tikzpicture}[scale=1.0]
  \def\wid{3.5}
  \def\ht{2}
  \node at (0,0 ) (bl) {};
  \node at (\wid,0) (br) {};
  \node at (\wid,\ht) (tr) {};
  \node at (0,\ht) (tl) {};
  \node at (0,\ht/2) (ml) {};
  \node at (\wid/3,\ht/2) (mml) {};
  \node at (2*\wid/3,\ht/2) (mmr) {};
  \node at (\wid,\ht/2) (mr) {};
  \node at (\wid/2,2*\ht/3) (mup) {};
  \node at (\wid/3,0) (bml) {};
  \node at (2*\wid/3,0) (bmr) {};
  \node at (\wid/2,\ht) (tm) {};
  \draw (bl.center)--(br.center)--(tr.center)--(tl.center)--(bl.center);
  \draw (ml.center)--(mml.center)--(mup.center)--(mmr.center)--(mr.center);
  \draw (bml.center)--(mml.center);
  \draw (bmr.center)--(mmr.center);
  \draw (mup.center)--(tm.center);
\end{tikzpicture}
\end{center}

We would like to know if it is possible to cross each interior wall in the building exactly once (without teleporting).

\begin{enumerate}[(a)]
  \item Turn this into graph theory problem. (Draw the corresponding graph.)
\item Either find such a walk, or prove that no such walk exists.
\item What if we want to pass through the exterior walls exactly once as well?
\end{enumerate}

\end{prob}

\begin{prob} 
  If~$G$ is a graph and~$e$ is an edge of~$G$, define~$G-e$ to be the graph
  obtained from~$G$ by removal of~$e$ (but not the endpoints of~$e$).
  Recall that to say a graph is {\em connected} means that every pair of its
  vertices can be connected by a path. 
  \begin{enumerate}[(a)]
    \item Give an example of a connected graph~$G$ with an edge~$e$ such
      that~$G-e$ is not connected.
    \item Suppose~$G$ is a connected graph and~$e$ is an edge of~$G$ that is
      part of a cycle.  Prove that removal of~$e$ does not disconnect the graph. 
      Your proof is required to start with the line: ``Let~$u$ and~$v$ be vertices
      of~$G$.'' It should then show there must be a path in~$G-e$ 
      connecting~$u$ and~$v$.
  \end{enumerate}
\end{prob}

\end{document}
